\documentstyle[% 


righttag% 


,11pt% 


,amssymb% 


,russian%               remove in Eng. 


,izvru%                 Set izven% in Eng. 


% ,izvru-ks             for brief communications 


]{amsart} 


 


%       Math definitions 


 


\newcommand{\rank}{\operatorname{rank}} 


\newcommand{\tts}[1]{} 


 


\theoremstyle{plain} 


\theorembodyfont{\it} 


\newtheorem{lemnonum}{\hskip\parindent Лемма} 


\renewcommand{\thelemnonum}{} 


 


 


%\hoffset=-15mm%                  For Eng. 


  


\setcounter{page}{32} 


\god{1998} 


\nomer{10~(437)} %                        Remove in Eng. 


%\nomer{Vol.\,42, No.\,10}%               Set in Eng. 


%\str{\pageref{begin}--\pageref{end}}%  Set in Eng. 


\udk{512.926} 


\author{\it В.П.\,Деревенский} 


% NB:   ^^ remove in Eng. 


\title{ Решение систем матричных алгебраических\\ 


линейных уравнений} 


 


\date{} 


% \pagestyle{myheadings}%         Set in Eng. 


 


\pagestyle{plain} %              Remove in Eng. 


\begin{document} 


 


% \thispagestyle{plain}%          Set in Eng. 


 


\maketitle 


\label{begin} 


 


Основным методом решения систем матричных алгебраических линейных 


левосторонних уравнений (СМАЛЛУ) 


\begin{equation} 


A^j_i X_j=B_i,\quad i,j=\ov{1,N}, 


\end{equation} 


где по повторяющимся верхним и нижним индексам производится суммирование, 


а $A^j_i$ и $B_i$ принадлежат $M_n$ (множеству квадратных матриц в $R^n$), 


как 


известно, является ``обобщенный алгоритм Гаусса'' (\cite{1}, с.\,55), к 


сожалению, не дающий в общем случае формул для выражения $X_j$ через 


данные матричные параметры. Ниже предлагается способ определения таких 


формул. 


 


Введем несколько понятий. 


 


\begin{defn} 


Левым оператором Гаусса невырожденного элемента матричной 


последовательности ($A_i$, $i=\ov{1,N}$) называется оператор 


$\Gamma(A_i)$, 


действующий на другую матричную последовательность $(Y_i)$ по правилу 


$$ 


\Gamma(A_i)Y_j=Y_j-A_j(A_i)^{-1}Y_i,\quad |A_i|\ne 0. 


$$ 


\end{defn} 


 


\begin{defn} 


Произведением левых операторов Гаусса называется 


$$ 


\Gamma(B_k)\Gamma(A_i)Y_j=\Gamma(\Gamma(A_i)B_k) 


(\Gamma(A_i)Y_j). 


$$ 


\end{defn} 


 


Так как 


$\Gamma(A^j_i)A^j_k=A^j_k-A^j_k(A^j_i)^{-1}A^j_i=0$, 


то введенный оператор осуществляет ``обобщенный алгоритм 


Гаусса'' решения СМАЛЛУ. 


 


Если матрица $A^1_1$ в системе (1) невырожденная (в противном случае 


производится перенумерация уравнений и неизвестных), то 


последовательное применение 


$\Gamma(A^1_1)$ ко всем уравнениям, начиная со второго, 


обращает в нуль в них все коэффициенты при $X_1$, приводя к виду 


$$ 


A^j_1X_j=B_1;\quad 


\Gamma(A^1_1)A^k_l X_k=\Gamma(A^1_1)B_l,\quad 


k,l=\ov{2,N}. 


$$ 


Если матрица $\Gamma(A^1_1)A^2_2$ 


обратима, то, преобразуя с 3-го по последнее уравнения


этой системы оператором $\Gamma(\Gamma(A^1_1)A^2_2)$ и т.\,д., 


получаем верхнетреугольную форму системы (1) 


\begin{equation} 


A^{*j}_i X_j=B^*_i,\quad j=\ov{i,N}, 


\end{equation} 


в которой в силу определения 2 матричные параметры преобразуются по 


закону 


$$ 


B^*_i=\prod^{i-1}_{k=1}\Gamma(A^k_k)B_i,\quad 


\prod^N_{i=1}A_i=A_N A_{N-1}\dotsc A_1. 


$$ 


Разумеется, согласно определению 1 триангуляция достижима лишь в случае 


невырожденности всех параметров операторов Гаусса. В связи с этим 


возникает необходимость ввести аналогичное скалярному случаю


 


\begin{defn} 


Система (1) называется невырожденной, если существует такая нумерация 


индексов $i$ и $j$, для которой матрица 


$$ 


D(A)=A^1_1\prod^{N-1}_{i=2}\biggl(\,\prod^1_{k=i-1} 


\Gamma(A^{i-k}_{i-k})A^i_i\biggr),\quad A=(A^j_i), 


$$ 


является невырожденной, 


\begin{equation} 


d(A)=|D(A)|\ne 0. 


\end{equation} 


\end{defn} 


 


Таким образом, если система (1) является невырожденной, то она 


триангулируется в форме (2) ``прямым ходом'' обобщенного алгоритма 


Гаусса. ``Обратный ход'', т.\,е.{} диагонализация системы, 


осуществляется введенным оператором Гаусса в обратном порядке с


использованием операторов 


$\Gamma(A^{*i}_i)$ ($i=\ov{N,2}$). Так как при этом не преобразуются 


диагональные матрицы полученной системы, то (2) диагонализируется в 


форме 


\begin{equation} 


A^{*i}_i X_i=\prod^{i+1}_{\tau=N}\Gamma(A^{*\tau}_\tau)B^*_i, 


\end{equation} 


что в силу условия (3) позволяет записать неизвестные матрицы в виде 


$$ 


X_i=(A^{*i}_i)^{-1}\prod^{i+1}_{\tau=N} 


\Gamma(A^{*\tau}_\tau)B^*_i. 


$$ 


Подставляя сюда $A^{*i}_i$, имеем 


\begin{equation} 


X_i=\bigg[\,\prod^{i-1}_{k=1}\Gamma(A^k_k)A^i_i\bigg]^{-1} 


\prod^{i+1}_{\tau=N}\Gamma(A^\tau_\tau)\bigg(\, 


\prod^{i-1}_{k=1}\Gamma(A^k_k)B_i\bigg). 


\end{equation} 


 


Следовательно, справедлива 


 


\begin{thm} 


Перенумерацией уравнений и неизвестных невырожденной системы $(1)$ 


можно привести ее решение к виду $(5)$. 


\end{thm} 


 


Предложенный метод решения матричных линейных систем с односторонним 


умножением может быть представлен в более компактной форме следующим 


образом. 


 


\begin{defn} 


Верхним триангулятором системы (1) называется оператор $T$, действующий 


на блочную $(N\times N)$-мерную матрицу $A=(A^j_i)$ как диагональная 


операторная матрица, диагональными элементами которой являются 


$$ 


T^i_i=\prod^1_{k=i-1}\Gamma(A^k_k),\quad T^1_1=E, 


$$ 


где $E$ --- тождественный оператор, а нижним триангулятором системы (1) 


--- аналогичный оператор $T^*$, диагональными элементами которого 


являются 


$$ 


T^{*i}_i=\prod^{i+1}_{k=N}\Gamma(A^k_k),\quad T^N_N=E. 


$$ 


\end{defn} 


 


Процедуры получения треугольной и диагональной форм (2) и (4)


системы (1) сводятся к последовательному ``умножению'' слева обеих 


сторон ее в блочно-матричной записи 


$$ 


AX=B,\quad A=(A^j_i),\quad B=(B_i),\quad X=(X_i) 


$$ 


на $T$ и $T^*$. ``Умножение'' на $T$ определило (2), а последующее 


действие оператора $T^*$ определило (4): $T^*TAX=T^*TB$. Следовательно, 


$$ 


X=(T^*TA)^{-1}(T^*TB). 


$$ 


Из сопоставления этой формулы с (5) видно, что 


$(T^*TA)^{-1}=((A^{*i}_i)^{-1})$, $TB=(B^*_i)$, 


$T^*TB=\Big(\prod\limits^{i+1}_{\tau=N}\Gamma(A^{*\tau}_\tau) 


B^*_i\Big)$. 


 


Так как определение 3 и теорема 1 формулируются с точностью до способа 


нумерации матриц $A^j_i$, то перестановка операторов $T$ и $T^*$, 


очевидно, допустима. Так что вполне естественно считать основной 


верхнетреугольную форму преобразованной системы (1), получаемую с 


помощью оператора $T$, как это исторически сложилось при использовании 


метода Гаусса решения скалярных систем. Однако диагонализацию можно 


проводить и в обратном порядке, преобразуя сначала систему (1) к 


нижнетреугольному виду оператором $T^*$, а затем оператором $T$ --- к 


диагональному виду $X=(TT^*A)^{-1}(TT^*B)$. 


 


Аналогичные рассуждения справедливы, разумеется, и для систем матричных 


алгебраических правосторонних линейных уравнений (СМАПЛУ) 


\begin{equation} 


X_i A^i_j=B_j,\quad i,j=\ov{1,N}, 


\end{equation} 


для формулировки которых необходимо несколько изменить введенные 


понятия. 


 


\begin{defn} 


Правым оператором Гаусса невырожденной матрицы $A_i$ из 


последовательности $(A_j)$ называется оператор $\Gamma(A_i)$, 


действующий на другую матричную 


последовательность $(Y_j)$ по правилу 


$$ 


Y_j\Gamma(A_i)=Y_j-Y_i(A_i)^{-1}A_j. 


$$ 


\end{defn} 


 


\begin{defn} 


Произведением правых операторов Гаусса называется 


\begin{equation} 


Y_j\Gamma(A_i)\Gamma(B_k)=(Y_j\Gamma(A_i)) 


\Gamma(B_k\Gamma(A_i)). 


\end{equation} 


\end{defn} 


 


\begin{defn} 


Система (6) называется невырожденной, если существует такая нумерация 


индексов $i$ и $j$, что 


\begin{equation} 


d\bigg(A^1_1\prod^{N-1}_{i=2}\bigg(A^i_i


\prod^{i-1}_{k=1}\Gamma(A^{i-k}_{i-k})\bigg)\bigg)\ne 0. 


\end{equation} 


\end{defn} 


 


С учетом этих переопределений аналог теоремы 1 будет формулироваться 


следующим образом. 


 


\begin{thm} 


Решение невырожденной системы $(6)$ при такой нумерации индексов $i$ и 


$j$, что выполняется условие $(8)$, имеет вид 


$$ 


X_i=B_i\prod^1_{k=i-1}\Gamma(A^k_k)\prod^N_{\tau=i+1} 


\Gamma(A^\tau_\tau)\bigg[A^i_i\prod^1_{k=i-1} 


\Gamma(A^k_k)\bigg]^{-1}, 


$$ 


где операторы Гаусса и их произведения определяются в $(7)$ и $(8)$. 


\end{thm} 


 


Следует отметить, что в скалярных случаях приведенные утверждения 


определяют обычные способы решения алгебраических систем с неизвестными 


векторами столбцов и строк соответственно. 


 


Рассмотрим более сложный случай систем матричных алгебраических 


линейных двусторонних уравнений (СМАЛДУ) 


\begin{equation} 


\sum^N_{j=1}\sum^{k(i,j)}_{k=1}A_{ijk}X_j B_{ijk}=C_i,


\quad (A_{ijk},B_{ijk},C_i)\subset M_n, 


\end{equation} 


для решения которых введем два новых определения. 


 


\begin{defn} 


Вектор (столбец) с компонентами $\wt X_u$ называется редуцированной 


матрицей $X=(x^q_p)\in M_n$ ($p,q=\ov{1,n}$), если 


\begin{equation} 


\wt X_u=\wt X_{n(p-1)+q}=x^q_p. 


\end{equation} 


\end{defn} 


 


\begin{defn} 


Тензорное (прямое) произведение матрицы $A$ на транспонированную 


матрицу 


$B$ называется их редуцентом и обозначается $R(A,B)$, 


$$ 


R(A,B)=A\otimes B'. 


$$ 


\end{defn} 


 


Для решения системы (9) редуцируем ее, записав в векторной форме в 


$R^{n^2}$. 


Для этого необходимо произвести перенумерацию (10) компонент всех 


матриц системы. При этом справедлива \cite{2} 


 


\begin{lemnonum} 


Двустороннее умножение матрицы $X$ на квадратные матрицы $A$ и $B$ 


эквивалентно левому умножению редуцированной матрицы $X$ на $R(A,B)$. 


\end{lemnonum} 


 


В соответствии с этим утверждением система (9) принимает вид 


\begin{equation} 


\sum^N_{j=1}R_{iju}^v\wt X_{jv}=\wt C_{iu},\quad u,v=\ov{1,n^2}, 


\end{equation} 


где $R_{ij}=\sum\limits^{k(i,j)}_{k=1}R(A_{ijk}, B_{ijk})$, т.\,е.{} 


принимает вид системы (1) с векторными неизвестными и 


свободными членами, компоненты которых отмечены вторыми индексами $u$ и 


$v$, пробегающими значения от 1 до $n^2$. 


 


Объединяя пары индексов $(i,u)$ и $(j,v)$ в единые индексы 


$\alpha=N(i-1)+u$ и $\beta=N(j-1)+v$, переходом из 


$R^{n^2}$ в $R^{Nn^2}$ систему (11) можно записать как обычную 


скалярную линейную неоднородную систему 


\begin{equation} 


R_\alpha^\beta\wt x_\beta=\wt c_\alpha,\quad 


\alpha,\beta=\ov{1,Nn^2}. 


\end{equation} 


 


Этот вид системы (9) позволяет утверждать следующее. 


 


\begin{thm} 


Для совместности системы $N$ матричных алгебраических линейных 


двусторонних уравнений с неизвестными $(9)$, необходимо и достаточно, 


чтобы


$$


\rank(R^\beta_\alpha)=\rank(R^\beta_\alpha,\wt c_\alpha). 


$$


\end{thm} 


 


При невырожденности системы (12) $\wt x_\alpha$ находятся любым 


классическим 


методом. Однако из-за большой размерности пространства $R^{Nn^2}$ имеет 


смысл обратиться к (11), чтобы воспользоваться предложенным выше 


методом, который действует и в случае искомых матриц, имеющих 


размерность $n\times 1$. Для этого введем 


 


\begin{defn} 


Система (9) называется невырожденной, если при некоторой перенумерации 


индексов $i$ и $j$\; $d(R)\ne 0$, $R=(R_{ij})$. 


\end{defn} 


 


Теперь на основании теоремы 1 может быть сформулирована 


 


\begin{thm} 


Решением невырожденной матричной системы $(9)$ являются матрицы $X_i$ в 


$R^n$, компоненты которых определяются в виде 


$$ 


x^q_{ip}=\wt x_{in(p-1)+q}, 


$$ 


где векторы $\wt x_i$ в $R^{n^2}$ с компонентами $\wt x_{in(p-1)+q}$ 


определяются как 


$$ 


\wt x_i=(R^*_{ii})^{-1}\prod^{i+1}_{\tau=N} 


\Gamma(R^*_{\tau\tau})c^*_i, 


$$ 


а $c^*_i=\prod\limits^{i-1}_{k=1}\Gamma(R_{kk})c_i$. 


\end{thm} 


 


\begin{thebibliography}{9} 


 


\bibitem{1} 


Гантмахер Ф.Р. {\em Теория матриц}. -- 3-е изд. -- М.: Наука, 1967. -- 


575 с. 


\bibitem{2} 


Деревенский В.П. {\em Матричные двусторонние линейные дифференциальные 


уравнения} // Матем. заметки. -- 1994. -- Т.\,55. -- \No\,1. -- 


С.\,35--42. 


 


\end{thebibliography} 


 


\vskip 2cm 


 


\post{\it Казанская архитектурно-}{\it Поступила} 


\post{\it строительная академия}{04.12.1995} 


 


\label{end} 


\end{document} 


